\begin{pastexam}[label={px:nccu-115-dsa-03}]{政大資科所 115 學年度 第 3 題}
（單選）

(5\%) In a binary search tree, if an element is searched for, and the search path visits nodes with
keys 50, 75, 60, 68. What is the relationship between the key 65 and the node with key 68?
\begin{enumerate}[label=(\Alph*)]
\item 65 must be in the left subtree of 68.
\item 65 must be in the right subtree of 68.
\item 65 could be either in the left or right subtree of 68.
\item 65 must be a child of 68.
\end{enumerate}
\end{pastexam}

\begin{pastsoln}
答案是 (A)。

題目的意思是：$65$ 也在這棵樹裡，問它的位置。

\textbf{從搜尋路徑讀出每一步往哪走}。 二元搜尋樹的搜尋（資結書虛擬碼~9.1）在每個節點比較一次：
要找的鍵比節點小就往左，比節點大就往右。
\begin{itemize}
\item $50 \to 75$：$75 > 50$，所以 $75$ 是 $50$ 的右兒子，往右走。
\item $75 \to 60$：$60 < 75$，$60$ 是 $75$ 的左兒子。
\item $60 \to 68$：$68 > 60$，$68$ 是 $60$ 的右兒子。
\end{itemize}
由二元搜尋樹的性質，$68$ 這棵子樹裡的鍵都比 $50$ 大（在 $50$ 的右子樹裡）、比 $75$ 小（在 $75$
的左子樹裡）、比 $60$ 大（在 $60$ 的右子樹裡），也就是都在開區間 $(60, 75)$ 裡。

\textbf{65 在哪裡}。 在這棵樹裡搜尋 $65$：$65 > 50$ 往右到 $75$；$65 < 75$ 往左到 $60$；
$65 > 60$ 往右到 $68$；$65 < 68$，往左。 搜尋一定會找到它，所以 $65$ 在 $68$ 的左子樹裡。

逐項檢查：
\begin{itemize}
\item (A) 對，理由如上。
\item (B) 錯。 $68$ 的右子樹裡的鍵都大於 $68$，$65$ 不可能在那裡。
\item (C) 錯。 只可能在左子樹。
\item (D) 錯。 不一定是兒子。 例如 $68$ 的左兒子是 $62$，$65$ 掛在 $62$ 的右邊，搜尋 $65$ 的路徑是
$50, 75, 60, 68, 62, 65$，$65$ 是 $68$ 的孫子。
\end{itemize}
\end{pastsoln}
